\documentclass[11pt,a4paper]{article}
\usepackage[italian]{babel}
\usepackage[margin=2.5cm]{geometry}
\usepackage{parskip}
\usepackage{amsmath, amssymb, amsthm}
\usepackage{booktabs}
\usepackage{array}
\usepackage{xcolor}
\usepackage{needspace}
\usepackage{enumitem}
\usepackage{listings}
\usepackage{tikz}
\usepackage[most]{tcolorbox}
\usepackage[hidelinks]{hyperref}
\usetikzlibrary{decorations.pathreplacing, shapes.geometric, arrows.meta, positioning}

% Niente vedove/orfane: i paragrafi non lasciano righe isolate a fine/inizio pagina
\widowpenalty=10000
\clubpenalty=10000

% Elenchi più compatti
\setlist{itemsep=2pt, parsep=0pt, topsep=4pt, partopsep=0pt}

% --- Riquadri con bordo nero, senza numero (si spezzano tra due pagine) ---
% Uso: \begin{definizione} ... \end{definizione}
%      \begin{teorema}[Nome] ... \end{teorema}  (il nome è facoltativo)
\tcbset{nota/.style={enhanced jigsaw, breakable, colback=white,
  colframe=black, boxrule=0.5pt, sharp corners}}
\NewTColorBox{definizione}{}{nota,
  before upper={\textbf{Definizione.}\ }}
\NewTColorBox{teorema}{o}{nota,
  before upper={\textbf{Teorema\IfValueT{#1}{ (#1)}.}\ }}
\NewTColorBox{proposizione}{o}{nota,
  before upper={\textbf{Proposizione\IfValueT{#1}{ (#1)}.}\ }}
\NewTColorBox{proprieta}{o}{nota,
  before upper={\textbf{Proprietà\IfValueT{#1}{ (#1)}.}\ }}
\NewTColorBox{assioma}{o}{nota, boxrule=1pt,
  before upper={\textbf{Assioma\IfValueT{#1}{ (#1)}.}\ }}

% --- Ambienti semplici (senza riquadro) ---
% Il titolo sta in cima al blocco, anche se il contenuto inizia con un riquadro o
% con codice; e non resta mai da solo in fondo a una pagina.
% Uso: \begin{esempio}[Titolo facoltativo] ... \end{esempio}
\newcommand{\newrunin}[2]{%
  \NewDocumentEnvironment{#1}{o}{%
    \par\Needspace{4\baselineskip}\addvspace{\medskipamount}\noindent
    \textbf{#2}\IfValueT{##1}{ (##1)}\textbf{.}\ \ignorespaces}%
  {\par\addvspace{\medskipamount}}}
\newrunin{esempio}{Esempio}
\newrunin{esercizio}{Esercizio}
\newrunin{osservazione}{Osservazione}

% V e F nelle tavole di verità
\newcommand{\V}{\textsf{V}}
\newcommand{\F}{\textsf{F}}

% --- Codice C ---
% Uso: \begin{codice} ... \end{codice}      codice C numerato
%      \begin{comando} ... \end{comando}    comandi da terminale
%      \begin{risultato} ... \end{risultato} output di un programma
%      \lstinline|printf("ciao");|           codice dentro al testo
\lstdefinestyle{c}{
  language=C,
  basicstyle=\ttfamily\small,
  keywordstyle=\color{blue}\bfseries,
  commentstyle=\color{gray}\itshape,
  stringstyle=\color{red!70!black},
  numbers=left,
  numberstyle=\tiny\color{gray},
  numbersep=8pt,
  columns=flexible,
  keepspaces=true,
  breaklines=true,
  showstringspaces=false,
  tabsize=4,
  morekeywords={include, define},
  % lettere accentate nei commenti
  literate={à}{{\`a}}1 {è}{{\`e}}1 {é}{{\'e}}1 {ì}{{\`i}}1
           {ò}{{\`o}}1 {ù}{{\`u}}1 {È}{{\`E}}1 {À}{{\`A}}1
}
\lstset{style=c}
\newtcblisting{codice}{listing only, nota, left=6mm,
  listing options={style=c}}
\newtcblisting{comando}{listing only, nota,
  listing options={basicstyle=\ttfamily\small, numbers=none, language={},
    columns=flexible, keepspaces=true, breaklines=true}}
\newtcblisting{risultato}{listing only, enhanced jigsaw, breakable,
  colback=black!5, colframe=black, boxrule=0.5pt, sharp corners,
  listing options={basicstyle=\ttfamily\small, numbers=none, language={},
    columns=flexible, keepspaces=true, breaklines=true}}

% --- Diagrammi di flusso (simboli standard) ---
\tikzset{
  terminale/.style = {ellipse, draw, minimum width=2.5cm, minimum height=0.9cm},
  io/.style        = {trapezium, trapezium left angle=70, trapezium right angle=110,
                      draw, minimum width=2.5cm, minimum height=0.9cm},
  processo/.style  = {rectangle, draw, minimum width=2.5cm, minimum height=0.9cm},
  decisione/.style = {diamond, draw, aspect=2, minimum width=2.5cm},
  freccia/.style   = {thick, -{Stealth}}
}

\title{Funzioni invertibili e funzioni monotone}
\author{Marco Cerbella}
\date{Lezione 5 -- 8 ottobre 2026}

\begin{document}
\maketitle

\section{Funzioni invertibili}

\begin{definizione}
$f : D \to \mathbb{R}$ si dice \emph{invertibile} se per ogni uscita $y \in f(D)$ esiste un solo ingresso $x \in D$ tale che $f(x) = y$. La \emph{funzione inversa} di $f$ è la funzione che associa a ogni $y \in f(D)$ l'unico $x \in D$ tale che $f(x) = y$; si indica con $f^{-1}$.
\end{definizione}

Significato geometrico: il grafico di $f$ è intersecato al massimo in un punto da ogni retta parallela all'asse delle ascisse.

\begin{itemize}
    \item funzione biunivoca $\Rightarrow$ funzione invertibile;
    \item funzione iniettiva $\Rightarrow$ funzione invertibile su $\mathrm{Im}(f)$.
\end{itemize}

Il grafico di $f^{-1}$ è
\[
\mathrm{Gr}(f^{-1}) = \{(y, f^{-1}(y)) : y \in D(f^{-1})\} = \{(f(x), x) : x \in D\},
\]
cioè si ottiene da quello di $f$ scambiando tra loro le componenti di ciascuna coppia. Geometricamente, il grafico di $f^{-1}$ si ottiene ribaltando il grafico di $f$ rispetto alla bisettrice del primo e terzo quadrante.

\begin{esempio}
\leavevmode
\begin{itemize}
    \item $f(x) = x + 1$ è invertibile: da $x + 1 = y$ si ricava $x = y - 1$, quindi $f^{-1}(x) = x - 1$. I punti $(0,1)$ e $(1,0)$ del grafico di $f$ e di $f^{-1}$ sono simmetrici rispetto alla bisettrice.
    \item $f(x) = ax + b$ è invertibile per $a \neq 0$ e $f^{-1}(x) = \dfrac{x - b}{a}$.
    \item $g(x) = x^2$ \emph{non} è invertibile su $\mathbb{R}$, ma lo è se si restringe il dominio:
    \begin{itemize}
        \item $g|_{[0,+\infty)}$ è invertibile e $(g|_{[0,+\infty)})^{-1}(x) = \sqrt{x}$;
        \item $g|_{(-\infty,0]}$ è invertibile e $(g|_{(-\infty,0]})^{-1}(x) = -\sqrt{x}$.
    \end{itemize}
    \item $h(x) = x^3$ è invertibile e $h^{-1}(x) = \sqrt[3]{x}$.
\end{itemize}
\end{esempio}

\section{Funzioni monotone}

\begin{definizione}
$f : D \to \mathbb{R}$ si dice
\begin{itemize}
    \item \emph{monotona crescente} se $\forall x_1, x_2 \in D$, $x_1 > x_2 \Rightarrow f(x_1) \geq f(x_2)$;
    \item \emph{monotona decrescente} se $\forall x_1, x_2 \in D$, $x_1 > x_2 \Rightarrow f(x_1) \leq f(x_2)$.
\end{itemize}
Se al posto di $\geq$ vale $>$ si dice \emph{strettamente crescente}; se al posto di $\leq$ vale $<$ si dice \emph{strettamente decrescente}.
\end{definizione}

\begin{proposizione}
Se $f : D \to \mathbb{R}$ è strettamente monotona allora è iniettiva, cioè invertibile su $\mathrm{Im}(f)$ (monotonia stretta $\Rightarrow$ invertibilità).
\end{proposizione}

Non vale il viceversa: invertibile $\not\Rightarrow$ monotona. Ad esempio
\[
f(x) = \begin{cases} \dfrac{1}{x} & \text{se } x \neq 0, \\[4pt] 0 & \text{se } x = 0 \end{cases}
\]
è invertibile ma non monotona: presi $x_2 < 0 < x_1$ si ha $x_1 > x_2$ e $f(x_1) > f(x_2)$, mentre per $0 < x_2 < x_1$ si ha $x_1 > x_2$ e $f(x_1) < f(x_2)$.

\begin{osservazione}
La somma di funzioni monotone concordi è monotona dello stesso tipo.
\end{osservazione}

\begin{esempio}
$f(x) = x^5 + x$ è monotona crescente, quindi invertibile, ma non si può esplicitare $f(x) = y$ nella forma $x = f^{-1}(y)$.
\end{esempio}

\begin{esempio}[Test]
Una funzione $f : \mathbb{R} \to [0,1]$ è sicuramente \emph{limitata} (risposta corretta: 1). Non è detto che sia negativa (ha valori in $[0,1]$), né iniettiva, né suriettiva: ad esempio $f(x) = 1$ non è suriettiva neppure considerando come codominio l'intervallo $[0,1]$.
\end{esempio}

\end{document}
